% !TeX root = ../drafts/04-committing-the-witness.tex
\section{Committing the witness}\label{sec:rs}

\subsection{Multilinear stacking}\label{sec:stacking}

Suppose we want to commit to $T$ multilinears $P_1,\dots,P_T$ over boolean hypercubes of various sizes. We could commit to each independently, at the cost of a large cumulative opening proof. Instead they are concatenated into a large multilinear polynomial $q$, which we commit to with an inner-product commitment scheme (Definition~\ref{def:ipcs}), with a protocol to reduce opening claims on $P_i$ to a claim on $q$. Order them by size, $\kappa_1\ge\dots\ge\kappa_T$ variables, and concatenate their values end to end,
\[
  q \;=\; P_1 \,\Vert\, P_2 \,\Vert\,\dots\,\Vert\, P_T \,\Vert\, 0,
\]
zero-padded to length $2^{M}$, with $M=\bigl\lceil\log_2\textstyle\sum_i 2^{\kappa_i}\bigr\rceil$. Placed largest first, $P_i$ starts at an offset that is a multiple of $2^{\kappa_i}$; writing $\mathsf{sel}_i\in\cube{M-\kappa_i}$ for the high bits of that offset,
\[
  \widetilde P_i(z)\;=\;\widetilde q(z,\mathsf{sel}_i)\qquad(z\in\cube{\kappa_i}).
\]
An evaluation claim $\widetilde P_i(\zeta)=c$ is therefore a claim on the stack,
\[
  \widetilde q(\zeta,\mathsf{sel}_i)\;=\;\sum_{w\in\cube{M}}\eq\bigl((\zeta,\mathsf{sel}_i),w\bigr)\,\widetilde q(w)\;=\;c,
\]
a weighted sum over $\widetilde q$ with weight $W(w)=\eq((\zeta,\mathsf{sel}_i),w)$.

\paragraph{Batching.} By the identity above, the claims the protocol raises on the $P_i$ transfer to a list of claims $\langle W_j,\widetilde q\rangle=c_j$ on the stack, with $\E$-valued weights and targets. The verifier then samples $\lambda\in\E$, and it remains to prove
\[
  \Bigl\langle \sum_j\lambda^{j-1}W_j,\;\widetilde q\Bigr\rangle\qeq\sum_j\lambda^{j-1}c_j
\]
via the PCS (Annex~\ref{annex:pcs}).

\paragraph{Padding-free commitment.} The columns fill $[0,P)$ with $P=\sum_i 2^{\kappa_i}$, and $q$ is zero from there to $2^{M}$, so up to half of what is committed is padding. Annex~\ref{annex:pcs} commits $q$ as an interleaved word of $2^{\ell_0}$ rows of $2^{M-\ell_0}$ values, and it is free to take the row index to be any $\ell_0$ coordinates of $q$; taking the TOP $\ell_0$ makes row $u$ the contiguous block $q\bigl[u\cdot 2^{M-\ell_0},\,(u+1)\cdot 2^{M-\ell_0}\bigr)$, so the padding is whole rows and only $\nrows=\bigl\lceil P/2^{M-\ell_0}\bigr\rceil$ of them carry data. The prover encodes those $\nrows$ rows alone; the rest are zero rows of the oracle, which it supplies as zeros wherever a column is read. A column still has its full $2^{\ell_0}$ entries, so the oracle the verifier reads is unchanged and this is a prover-side saving of $1-\nrows 2^{-\ell_0}$ on the encoding, bounded by one half because $M=\lceil\log_2 P\rceil$ keeps $\nrows>2^{\ell_0-1}$. The Merkle realization of \S\ref{sec:e2e} takes two further savings from the same structure by reading a column from the TOP row index downwards, so that the absent rows are the leading words of the hashed leaf: their whole compression blocks are one chaining value shared by every leaf, computed once rather than per leaf, and only the rest of the leaf travels in the proof, the verifier restoring the zeros it knows are there. Both are quantised to whole blocks, and the verifier derives $\nrows$ from the announced sizes to read a leaf. Two things make it sound, both discharged above. Every weight of \S\ref{sec:stacking} is supported on a placed column, hence vanishes on $[P,2^{M})$, which is what lets the omitted rows hold anything (Annex~\ref{annex:pcs}, ``Rows the prover may omit''); an out-of-domain weight would not vanish there, which is why level $0$ takes no OOD sample. And the fold now binds $q$'s top $\ell_0$ coordinates first, so the fold challenges arrive in an order rotated from $q$'s own: reading them in the order the levels draw them, the first $\ell_0$ bind $q$'s top $\ell_0$ coordinates and the remainder bind $q$'s coordinates $0,\dots,M-\ell_0-1$ in order. A weight written in $q$'s coordinates, as every weight in this section is, is therefore evaluated at that point rotated left by $\ell_0$, and a verifier applies that rotation once, immediately before evaluating it (\S\ref{sec:e2e}).

\subsection{Committing to booleans via Ring-Switching}\label{sec:ringswitch}

Flock's BLAKE2s argument (\S\ref{sec:tab-blake2s}) requires a commitment to a Boolean multilinear in $\nflock+\kskip$ variables. Ring switching~\cite{DP26} lets the prover pack its low $\kskip=6$ coordinates into the $\K$-valued multilinear $\qflock$ in $\nflock$ variables (Annex~\ref{annex:rs}), which occupies one region of the stack in \S\ref{sec:stacking}.
